%%%%%%%%%%%%%%%%%%%%%%%%%%%%%%%%%%%%%%%%%%%%%%%%%%%%%%%%%%%%%%%%%%%%%%
% Lớp 10 | Đại số tổ hợp | Quy tắc cộng-quy tắc nhân
% Dạng 0D8?1-5: Bài toán đếm số tự nhiên
% Số câu: 18  (sinh tự động bởi tach_cau.py; từ nay thầy sửa trực tiếp file này)
%%%%%%%%%%%%%%%%%%%%%%%%%%%%%%%%%%%%%%%%%%%%%%%%%%%%%%%%%%%%%%%%%%%%%%

%%%%%%%%%%%%%%%% TRẮC NGHIỆM 4 PHƯƠNG ÁN %%%%%%%%%%%%%%%%

% nguồn: MaHoaNH_TracNghiem_Lop10
\begin{ex}%[0D8H1-5]%[ma:4fcdcae7]
	Từ $X=\left\{1,2,3,4,5,6,7,8,9\right\}$ có bao  nhiêu cách chọn $1$ số hoặc chẵn hoặc là nguyên tố?
	 \choice
		{$4$} 
		{$5$} 
		 {$6$}
		{\True $7$} 
	\loigiai{ 
		Gọi $A$ là tập  con của $X$ chứa số chẵn. Khi đó,
		\[A=\left\{2;4;6;8\right\}\Rightarrow  n\left(A\right)=4.\]
		Gọi $B$ là tập con của $X$ mà chứa số  nguyên tố. Khi đó,
		\[B=\left\{2;3;5;7\right\}\Rightarrow n\left(B\right)=4.\] 
		Ta có $A\cap B=\{2\}\Rightarrow n\left(A\cup B\right)=n\left(A\right)+n\left(B\right)-n\left(A\cap  B\right)=4+4-1=7.$
	} 
\end{ex}

% nguồn: MaHoaNH_TracNghiem_Lop10
\begin{ex}%[0D8V1-5]%[ma:0981545e]
	Từ các chữ số $1$, $2$, $3$, $4$, $5$,  $6$ có thể lập được bao nhiêu số tự nhiên lẻ có $6$ chữ số khác nhau và trong mỗi số đó tổng của ba chữ số đầu lớn hơn tổng của ba chữ số cuối một đơn vị.
	\choice 
		 {$32$}
		{\True $72$} 
		 {$36$}
		 {$24$}
	\loigiai{ 
		Gọi $\overline{a_1 a_2 a_3 a_4 a_5 a_6}$  là số cần tìm.\\
		Vì số cần tìm là số lẻ nên $a_6 \in \left\{1; 3; 5\right\}$ và tổng ba chữ số  đầu lớn hơn tổng ba chữ số sau nên $\left(a_1+a_2+a_3 \right)-\left(a_4+a_5+a_6 \right)=1$.
		\begin{itemize}
			\item Với $a_6=1$ thì  $\left(a_1+a_2+a_3\right)-\left(a_4+a_5\right)=2$.\\
			Suy ra $\heva{&a_1, a_2, a_3\in \left\{2,3, 6\right\} \\&a_4, a_5\in \left\{4,5\right\}}$  hoặc $\heva{&a_1, a_2, a_3\in \left\{2,4, 5\right\} \\&a_4, a_5\in \left\{3,6\right\}.}$
			\item Với  $a_6=3$ thì $\left(a_1+a_2+a_3\right)-\left(a_4+a_5\right)=4$.\\
			Suy ra $\heva{&a_1, a_2, a_3\in \left\{2;  4; 5\right\} \\&a_4, a_5\in \left\{1,6\right\}}$ hoặc $\heva{&a_1, a_2, a_3\in \left\{1,4, 6\right\} \\&a_4, a_5\in \left\{2,5\right\}.}$
			\item  Với $a_6=5$ thì $\left(a_1+a_2+a_3\right)-\left(a_4+a_5\right)=6$.\\
			Suy ra $\heva{&a_1, a_2, a_3\in \left\{2,3, 6\right\} \\&a_4, a_5\in \left\{1,4\right\}}$ hoặc $\heva{&a_1, a_2,  a_3\in \left\{1,4, 6\right\} \\&a_4, a_5\in \left\{2,3\right\}.}$
		\end{itemize}
		Mỗi trường hợp có $3!\cdot2!=12$  số thỏa mãn yêu cầu.\\
		Vậy có tất cả $6\cdot12=72$ số  cần tìm.
	} 
\end{ex}


%%%%%%%%%%%%%%%% ĐÚNG/SAI %%%%%%%%%%%%%%%%

% nguồn: MaHoaNH_DungSai_Lop10
\begin{ex}%[0D8N1-5]%[ma:ad44ff23]
	Từ  các chữ số $1$, $2$, $3$, $4$, $5$, $6$.
	 \choiceTF
		{\True Có $720$ số tự nhiên có năm chữ số đôi một  khác nhau}
		{\True Có $120$ số tự nhiên có ba  chữ số đôi một khác nhau}
		{\True  Có $180$ số tự nhiên chẵn có bốn chữ số đôi một khác nhau}
		{Có $120$ số tự nhiên lẻ có bốn chữ số đôi  một khác nhau}
	\loigiai{ 
		Theo quy tắc  nhân, ta có $6\cdot 5\cdot 4\cdot 3\cdot 2=720$ số tự nhiên có năm chữ số đôi một khác nhau.\\
		Có $6\cdot 5\cdot 4=120$ số  tự nhiên có ba chữ số đôi một khác nhau.\\
		Có $3\cdot  5\cdot 4\cdot 3=180$ số tự nhiên chẵn có bốn chữ số đôi một khác nhau.\\
		Tương tự, có $3\cdot 5\cdot 4\cdot 3=180$  số tự nhiên lẻ có bốn chữ số đôi một khác nhau.
	} 
\end{ex}

% nguồn: MaHoaNH_DungSai_Lop10
\begin{ex}%[0D8H1-5]%[ma:e67c7a92]
	Cho hai mệnh đề $P$: \lq\lq $ABC$ là tam giác đều\rq\rq\ và mệnh đề $Q$: \lq\lq $ABC$ là  tam giác cân\rq\rq.
	Xét mệnh đề kéo theo \lq\lq Nếu $P$  thì $Q$\rq\rq.
	\choiceTF 
		{\True $P$ là điều kiện  đủ để có $Q$}
		{\True $Q$ là điều kiện cần để  có $P$}
		{\True Mệnh đề  \lq\lq Nếu $P$ thì $Q$\rq\rq\ là mệnh đề đúng}
		 {Mệnh đề \lq\lq Nếu $Q$ thì $P$\rq\rq\ là mệnh đề đúng}
	\loigiai{ 
		Mệnh  đề $P \Rightarrow Q$ được phát biểu: \lq \lq Nếu $ABC$ là tam giác đều thì $ABC$ là tam giác cân\rq \rq.
		\begin{itemchoice}
			\itemch  $P$ là điều kiện đủ để có $Q$.
			\itemch $Q$ là điều kiện cần để có $P$. 
			\itemch Mệnh đề \lq\lq Nếu $P$ thì $Q$\rq\rq\ là mệnh đề đúng vì mọi tam giác đều là tam giác cân. 
			\itemch Không phải  tam giác cân nào cũng là tam giác đều nên mệnh đề \lq\lq Nếu $Q$ thì $P$\rq\rq\ là mệnh đề sai.
		\end{itemchoice}
	} 
\end{ex}

% nguồn: MaHoaNH_DungSai_Lop10
\begin{ex}%[0D8H1-5]%[ma:2dd64443]
	Cho hai mệnh đề sau  \\
	$P$: \lq\lq Hình bình hành có một góc vuông là hình chữ  nhật\rq\rq.\\
	$Q$: \lq\lq Số $7$ là  hợp số\rq\rq.
	\choiceTF 
		 {\True Mệnh đề $P$ là mệnh đề đúng}
		{Mệnh đề $Q$  là mệnh đề đúng}
		{Mệnh đề $P \Rightarrow Q$ là mệnh đề đúng} 
		{Mệnh đề $Q \Rightarrow P$ là mệnh đề sai} 
	\loigiai{ 
		\begin{itemchoice}
			\itemch Mệnh đề $P$ là mệnh đề  đúng vì hình bình hành có một góc vuông chính là hình chữ nhật.
			\itemch Số $7$ là số nguyên  tố, không phải hợp số. Suy ra mệnh đề $Q$ là sai.
			\itemch $P$ đúng và $Q$ sai nên mệnh đề  kéo theo $P \Rightarrow Q$ là sai.
			\itemch $Q$ sai và $P$ đúng nên $Q \Rightarrow P$ là mệnh đề đúng. Do đó phát  biểu sai.
		\end{itemchoice}
	} 
\end{ex}

% nguồn: MaHoaNH_DungSai_Lop10
\begin{ex}%[0D8H1-5]%[ma:babe9581]
	Từ sáu chữ số $0$, $1$, $2$, $3$, $4$,  $5$ có thể lập được
	 \choiceTF
		{\True $600$ số tự nhiên có  sáu chữ số đôi một khác nhau}
		{$720$ số tự nhiên có có sáu chữ  số đôi một khác nhau}
		{$120$ số tự  nhiên có có sáu chữ số đôi một khác nhau}
		{\True $288$  số tự nhiên lẻ có sáu chữ số đôi một khác nhau}
	 \loigiai{
		Theo quy tắc nhân, ta có $5\cdot 5\cdot 4\cdot 3\cdot 2\cdot \cdot 1=600$ số tự nhiên có sáu chữ  số đôi một khác nhau.\\
		Gọi $\overline{abcdef}$ là số tự  nhiên lẻ có sáu chữ số đôi một khác nhau.\\
		 Chọn $f$ $3$ cách chọn.\\
		Chọn  $a$ $4$ cách chọn.\\
		Chọn $b$  có $4$ cách chọn.\\
		Chọn  $c$ có $3$ cách chọn.\\
		Chọn $d$ có  $2$ cách chọn.\\
		Chọn $e$ có $1$ cách  chọn.\\
		Vậy có  $3\cdot 4\cdot 4\cdot 3\cdot 2\cdot 1=288$ số tự nhiên lẻ có sáu chữ số đôi một khác nhau.
	} 
\end{ex}

% nguồn: MaHoaNH_DungSai_Lop10
\begin{ex}%[0D8H1-5]%[ma:98acb4b3]
	 Xét tính đúng, sai của các mệnh đề sau:
	 \choiceTF
		{$\forall n \in \mathbb{N},\ n^2$ chia  hết cho $7 \Rightarrow n$ chia hết cho $7$}
		{\True $\forall n \in \mathbb{N},\ n^2$ chia hết cho $5 \Rightarrow  n$ chia hết cho $5$}
		{\True Nếu tam giác $ABC$ không phải là tam giác đều thì tam giác đó có  ít nhất một góc nhỏ hơn $60^\circ$}
		{\True Tứ giác $ABCD$  nội tiếp được đường tròn $\Rightarrow\widehat{A} + \widehat{C}$ bằng $180^\circ$}
	\loigiai{ 
		\begin{itemchoice}
			\itemch Ta chứng minh bằng phương  pháp phản chứng. \\
			Giả sử $n$ không chia hết cho $7 \Rightarrow n = 7m + i$,  $i \in \{1,2,3,4,5,6\}$. \\
			Khi đó $n^2 = 49m^2 + 14im + i^2$, và  các giá trị $i^2$ đều không chia hết cho $7$ nên $n^2$ cũng không chia hết cho $7$.
			\itemch  Chứng minh tương tự như câu a.
			\itemch Giả sử tam giác $ABC$ không đều và mọi góc đều lớn hơn hoặc bằng  $60^\circ$.\\
			$\Rightarrow$ Tổng các góc $\geq 180^\circ$ và chỉ xảy ra khi cả ba góc bằng $60^\circ$, mâu thuẫn  với giả thiết. \\
			Do đó, tam giác $ABC$ phải có góc nhỏ hơn  $60^\circ$.
			\itemch Tứ giác nội  tiếp được đường tròn $\Rightarrow \widehat{A} + \widehat{C} = 180^\circ$.
		\end{itemchoice}
	 }
\end{ex}

% nguồn: MaHoaNH_DungSai_Lop10
\begin{ex}%[0D8H1-5]%[ma:f4dd9839]
	Cho  tập $A=\left\{ 1;2;3;4 \right\}$.
	\choiceTF 
		 {\True Có thể lập được $16$ số có $2$ chữ số từ các chữ số ở tập $A$}
		{Có thể lập được $16$ số có  $2$ chữ số khác nhau từ các chữ số ở tập $A$}
		{Có thể lập được  $8$ số chẵn có $2$ chữ số khác nhau từ các chữ số ở tập $A$}
		{\True Có thể lập được $8$ số lẻ có $2$ chữ số từ  các chữ số ở tập $A$}
	 \loigiai{
		\begin{itemchoice}
			 \itemch Gọi số cần tìm có dạng $\overline{ab}$ với $a;b\in A$.\\
			Vì  số cần tìm có $2$ chữ số nên $a$ có $4$ cách chọn, $b$ có $4$ cách chọn.\\
			Như vậy, ta có  $4\cdot4=16$ số có hai chữ số được lập từ tập hợp $A$.
			\itemch Gọi  số cần tìm có dạng $\overline{ab}$ với $a;b\in A$.\\
			Vì số cần tìm có $2$ chữ số khác nhau nên $a$ có $4$ cách chọn, $b$ có $3$ cách chọn.\\ 
			Như vậy, ta có $4\cdot3=12$ số có  hai chữ số khác nhau được lập từ tập hợp $A$.
			\itemch Gọi số cần tìm có dạng $\overline{ab}$  với $a;b\in A$.\\
			Vì số cần tìm là số chẵn có $2$ chữ số khác nhau nên $b$ có $2$ cách chọn, $a$ có $3$ cách  chọn.\\
			Như vậy, ta có $2\cdot3=6$ số chẵn có hai  chữ số khác nhau được lập từ tập hợp $A$.
			\itemch Gọi số cần tìm có dạng $\overline{ab}$  với $a;b\in A$.\\
			Vì số cần tìm là số lẻ có $2$ chữ số nên $b$ có $2$  cách chọn, $a$ có $4$ cách chọn.\\
			Như vậy, ta có $2\cdot4=8$  số lẻ có hai chữ số được lập từ tập hợp $A$.
		\end{itemchoice}
	 }
\end{ex}

% nguồn: MaHoaNH_DungSai_Lop10
\begin{ex}%[0D8H1-5]%[ma:e997d4c9]
	Cho các chữ số $1$, $2$, $3$. Xét tính  đúng sai của các khẳng định sau
	 \choiceTF
		{Từ các chữ số đã cho lập được $27$ số  có $3$ chữ số}
		{\True Từ các chữ số  đã cho lập được $9$ số có $3$ chữ số và là số chẵn}
		{\True Tổng các số có $3$ chữ số đôi một khác nhau  được lập từ các chữ số đã cho là $1332$}
		{Số các số $10$ chữ số được tạo thành từ các chữ số $1$, $2$, $3$ sao cho bất kì $2$ chữ số nào đứng cạnh nhau cũng hơn kém nhau $1$  đơn vị là $62$}
	\loigiai{ 
		\begin{itemchoice}
			\itemch  Từ các chữ số đã cho lập được $3\cdot3\cdot3 = 27$ số có $3$ chữ số.
			\itemch Số  cần lập có dạng $\overline{ab2}$. Số các số thoả mãn là $3\cdot3\cdot1 = 9$ số.
			\itemch Số các số có ba chữ số  đôi một khác nhau được lập từ các chữ số đã cho là $3\cdot2\cdot1 = 6$ số.\\
			Ở $6$ số này thì chữ số $1$, $2$, $3$ đều xuất hiện  $2$ lần ở các hàng đơn vị hay chục hay hàng trăm.\\
			Nên tổng là $\left(3+3+2+2+1+1 \right)\cdot\left(  100+10+1 \right)=1332$.\\
			\textbf{Chú ý:}  có thể cộng sáu số lập được là $123+132+231+213+312+321=1332$.
			\itemch Gọi số  tự nhiên cần tìm có dạng $\overline{a_1a_2\dots a_{10}}$.
			\begin{enumerate} [$\bullet$]
				\item \textbf{Bước 1:} Xếp chữ số $2$ ở vị trí lẻ $a_1$, $a_3$,$\dots$, $a_9$ hoặc vị trí chẵn $a_2$,  $a_4$,$\dots$, $a_{10}$ có $2$ cách.
				\item \textbf{Bước 2:} Xếp các chữ số  $1$ hoặc $3$ vào các vị trí còn lại có $2^5$ cách.\\
				Vậy  có $2^6=64$ số.
			\end{enumerate}
		\end{itemchoice}
	 }
\end{ex}

% nguồn: MaHoaNH_DungSai_Lop10
\begin{ex}%[0D8H1-5]%[ma:5b1e711e]
	Cho tập hợp $A=\left\{ 1,2,3,4,5 \right\}$. 
	 \choiceTF
		{\True Từ $A$ lập được $25$ số  có hai chữ số}
		{Từ $A$ lập được  $125$ số có ba chữ số khác nhau}
		{\True Từ $A$ lập được $24$ số chẵn có ba chữ số khác  nhau}
		{Từ $A$ lập được $101$ số lẻ có ba chữ số khác  nhau}
	\loigiai{ 
		\begin{itemchoice}
			\itemch  Gọi $\overline{ab}$ là số cần lập, số cách chọn $a$ là $5$, số cách $b$ là $5$.\\ Theo quy tắc nhân có $5\cdot5=25$ số có hai chữ số.
			\itemch Gọi số có $3$ chữ  số khác nhau là $\overline{abc}$.\\
			Chọn $a$ có  $5$ cách.\\
			Chọn $b$ có $4$  cách.\\
			Chọn  $c$ có $3$ cách.\\
			Suy  ra có $5\cdot4\cdot3=60$ số có ba chữ số khác nhau.
			\itemch Gọi số chẵn  có ba chữ số khác nhau là $\overline{abc}$.\\
			 Chọn $c$ có $2$ cách.\\
			Chọn  $a$ có $4$ cách.\\
			 Chọn $b$ có $3$ cách.\\
			Suy ra có $2\cdot4\cdot3=24$ số chẵn  có ba chữ số khác nhau.
			\itemch Gọi số lẻ có ba chữ số  khác nhau là $\overline{abc}$.\\
			Chọn $c$ có  $3$ cách.\\
			Chọn  $a$ có $4$ cách.\\
			Chọn $b$ có $3$  cách.\\
			Suy ra có $3\cdot4\cdot3=36$ số lẻ  có ba chữ số khác nhau.
		\end{itemchoice}
	} 
\end{ex}

% nguồn: MaHoaNH_DungSai_Lop10
\begin{ex}%[0D8H1-5]%[ma:45852b9c]
	Từ các chữ số $1$, $2$,  $3$, $4$, $5$. Xét tính đúng sai của các khẳng định sau
	 \choiceTF
		{\True Có thể lập được $5$ số có  một chữ số}
		{Có thể lập được $20$ số có  hai chữ số}
		{\True Có thể lập được $60$ số có ba chữ  số khác nhau}
		{Có thể lập được $32$ số có ba chữ số khác nhau không nhỏ  hơn $342$}
	\loigiai{ 
		\begin{itemchoice}
			\itemch Các số có một chữ số là $1,2,3,4,5$. 
			\itemch Gọi số cần tìm có dạng  $\overline{ab}$ trong đó $a,b\in \left\{ 1,2,3,4,5 \right\}$.\\
			Có $5$ cách chọn $a\in \left\{ 1,2,3,4,5 \right\}$ và có $5$ cách chọn $b\in  \left\{ 1,2,3,4,5 \right\}$.\\
			Theo quy tắc nhân, ta có  số các số cần tìm là $5.5=25$ (số).
			\itemch Gọi số cần tìm  có dạng $\overline{abc}$ trong đó $a,b,c$ khác nhau và $a,b,c\in \left\{ 1,2,3,4,5 \right\}$.\\
			Có $5$ cách chọn $a\in \left\{  1,2,3,4,5 \right\}$, có $4$ cách chọn $b$ (khác $a$) và có $3$ cách chọn $c$ (khác $a$ và khác $b$).\\
			Theo  quy tắc nhân, ta có số các số cần tìm là $5\cdot4\cdot3=60$ (số).
			\itemch Gọi số cần  tìm có dạng $\overline{abc}$ trong đó $a,b,c$ khác nhau và $a,b,c\in \left\{ 1,2,3,4,5 \right\}$.\\
			Từ giả thiết có $\overline{abc}\ge 342$, suy ra $a\ge 3$. Do đó,  $a\in \left\{ 3,4,5 \right\}$.
			\begin{enumerate} [$\bullet$]
				\item \textbf{TH1:} Nếu $a=3$ thì $\overline{3bc}\ge 342$ khi $b\ge 4$, tức  là $b\in \left\{ 4,5 \right\}$.\\
				Với $b=4$,  để $\overline{34c}\ge 342$ thì $c\ge 2$, $c$ khác $3$ và $4$ (vì ba chữ số khác nhau), suy ra có $2$ cách chọn $c\in \left\{ 2,5 \right\}$.\\
				Theo quy tắc nhân  có số các số có dạng $\overline{34c}$ là $1\cdot1\cdot2=2$ (số).\\
				Với $b=5$, để $\overline{35c}\ge 342$ thì có $3$ cách chọn $c\in \left\{ 1,2,4 \right\}$ (vì $c$  khác $3$ và $5$).\\
				Theo quy tắc nhân có số  các số có dạng $\overline{35c}$ là $1\cdot1\cdot3=3$ (số).\\
				Vậy theo quy tắc cộng, ta có số các số trong \textbf{TH1} là $2+3=5$  (số).
				\item \textbf{TH2:} Nếu $a=4$ thì $\overline{4bc}\ge  342$.\\
				Ta  có có $4$ cách chọn $b\in \left\{ 1,2,3,5 \right\}$ và có $3$ cách chọn $c\in \left\{ 1,2,3,5 \right\}\backslash \left\{ b \right\}$.\\
				Vậy theo quy tắc nhân có số các số có dạng $\overline{4bc}$ là  $1\cdot4\cdot3=12$ (số).
				 \item \textbf{TH3:} Nếu $a=5$ thì tương tự cũng có $12$ số dạng $\overline{5bc}$.
			\end{enumerate}
			Vậy theo quy tắc cộng, ta có tất cả  $5+12+12=29$ (số).
		\end{itemchoice}
	 }
\end{ex}

% nguồn: MaHoaNH_DungSai_Lop10
\begin{ex}%[0D8V1-5]%[ma:0c1137dd]
	Từ  các chữ số $0$, $1$, $2$, $3$, $4$, $5$, $6$, $7$.
	 \choiceTF
		{\True Có thể lập được $8$ số  có một chữ số}
		{Có thể lập được $56$ số có hai chữ số  khác nhau}
		{\True Có thể lập được $25$  số chẵn có hai chữ số khác nhau}
		{Có thể lập được  $22$ số có ba chữ số khác nhau chứa chữ số $2$ và chia hết cho $5$}
	 \loigiai{
		\begin{itemchoice}
			\itemch Có thể lập được $8$ số có một chữ  số.
			\itemch Gọi số cần tìm có  dạng $\overline{ab}$ trong đó $a\ne 0$, $b\in \left\{ 0,\,1,\,2,3,\,4,\,5,\,6,\,7 \right\}$.\\
			Có $7$ cách chọn  $a$ và mỗi cách chọn $a$ thì có $7$ cách chọn $b$. Vậy theo quy tắc nhân có $49$ số.
			\itemch Gọi số cần tìm có dạng $\overline{ab}$ trong đó $a$, $b\in \left\{ 0,\,1,\,2,3,\,4,\,5,\,6,\,7 \right\}$.\\ 
			TH1: số có dạng  $\overline{a0}$ có $7$ số vì $a$ có $7$ cách chọn.\\
			TH2: số có dạng $\overline{a2}$  có $6$ số vì $a$ có $6$ cách chọn.\\
			TH3: số có dạng $\overline{a4}$ có  $6$ số vì $a$ có $6$ cách chọn.\\
			TH4: số  có dạng $\overline{a6}$ có $6$ số vì $a$ có $6$ cách chọn.\\
			Vậy theo quy tắc cộng  có $25$ số chẵn có hai chữ số khác nhau.
			\itemch Gọi số cần tìm có dạng $\overline{abc}$ trong đó $a$, $b$, $c\in \left\{ 0,\,1,\,2,3,\,4,\,5,\,6,\,7 \right\}$.\\ 
			TH1: số có dạng $\overline{2b0}$ có $6$ số  vì $b$ có $6$ cách chọn.\\
			TH2: số có dạng $\overline{a20}$ có $6$ số  vì $a$ có $6$ cách chọn.\\
			TH3: số có dạng  $\overline{2b5}$ có $6$ số vì $b$ có $6$ cách chọn.\\
			TH4: số có dạng $\overline{a25}$ có  $5$ số vì $a$ có $5$ cách chọn.\\
			Vậy theo quy tắc cộng có $23$ số chẵn và chia hết cho $5$ có hai chữ số  khác nhau.
		\end{itemchoice}
	} 
\end{ex}

% nguồn: MaHoaNH_DungSai_Lop10
\begin{ex}%[0D8V1-5]%[ma:a5ecf17a]
	 Cho các chữ số $0,1,2,3,4,5,6,7$.
	 \choiceTF
		{\True Có thể lập được $8$ số tự nhiên có một chữ số} 
		{Có thể lập được  $56$ số tự nhiên có hai chữ số khác nhau}
		{\True Có thể lập được $25$ số tự nhiên chẵn có  hai chữ số khác nhau}
		{Có thể lập được $22$ số tự nhiên có ba chữ số khác nhau đôi một, chứa chữ  số $2$ và chia hết cho $5$}
	\loigiai{ 
		\begin{itemchoice}
			\itemch \textbf{Đúng}. 
			\itemch  \textbf{Sai}. \\
			Gọi $\overline{ab}$ là số  cần tìm. \\
			Chọn $a$ có $7$ cách và $b$ có $7$ cách.  \\
			Vậy  có $7\cdot7=49$ số.
			\itemch \textbf{Sai}.\\ 
			Gọi $\overline{ab}$ là số cần  tìm.
			\begin{itemize}
				\item TH1: $b=0$ có $1$ cách; chọn $b$ có $7$ cách.  Do đó có: $1\cdot 7=7$ số.
				\item TH2: $b\in\{2,4,6\}$: Chọn $b$ có $3$  cách và chọn $a$ có $6$ cách. do đó có $3\cdot 6=18$ số.
			\end{itemize}
			Vậy có $7+18=25$  số thỏa yêu cầu bài toán.
			\itemch \textbf{Sai}.\\ 
			Gọi $\overline{abc}$ là số  cần tìm.
			\begin{itemize}
				\item TH1: $a=2$: Có $2$ cách chọn $c$ và $6$ cách chọn $b$. Do đó có $2\cdot 6=12$  số.
				\item TH2: $b=2$, $c=5$: Chọn $a$ có $5$  cách.
				\item TH3:  $b=2$, $c=0$: Chọn $a$ có $6$ cách.
			\end{itemize}
			 Vậy có $12+5+6=23$ số.
		\end{itemchoice}
	} 
\end{ex}

% nguồn: MaHoaNH_DungSai_Lop10
\begin{ex}%[0D8V1-5]%[ma:652664d0]
	Lập các số tự nhiên từ các chữ số trong tập hợp  $A=\{0;1;2;3;4;5;6;7\}$. Các mệnh đề sau đúng hay sai?
	 \choiceTF
		{Lập được $336$ số  có ba chữ số khác nhau}
		{\True Lập được $144$  số là số lẻ có ba chữ số khác nhau}
		 {Lập được $84$ số có ba chữ số khác nhau và chia hết cho $5$}
		{\True Lập được $154$ số có ba chữ số khác nhau và nhỏ hơn $456$} 
	 \loigiai{
		\begin{itemchoice}
			\itemch Gọi số cần tìm có dạng $\overline{abc}$, với $a$, $b$, $c$ khác nhau lấy từ $A=\{0;1;2;3;4;5;6;7\}$ và $a \neq 0$. \\ 
			Số cách chọn $a$, $b$, $c$ khác nhau (theo thứ tự  đó) là $7$, $7$, $6$.\\
			Vậy theo quy tắc  nhân có $7 \cdot 7 \cdot 6=294$ số.
			\itemch Gọi số cần tìm  có dạng $\overline{abc}$ là số lẻ, với $a$, $b$, $c$ khác nhau lấy từ $A=\{0;1;2;3;4;5;6;7\}$ và $a \neq 0$. \\
			Chọn  $c \in \{1;3;5;7\}$ nên $c$ có $4$ cách chọn.\\
			Số cách chọn $a$, $b$ khác nhau (theo thứ tự đó)  là $6$, $6$.\\
			Vậy theo quy tắc nhân  có $4 \cdot 6 \cdot 6=144$ số.
			\itemch Gọi số cần tìm có dạng $\overline{abc} \;\vdots\; 5$, với $a,b,c$ khác nhau lấy từ $A=\{0;1;2;3;4;5;6;7\}$. Khi  đó $d=0$ hoặc $d=5$.
			\begin{itemize}
				\item Nếu $d=0$ ta xét số $\overline{ab0}$, với $a,b \in \{1;2;3;4;5;6;7\}$.\\ 
				Số cách chọn $a,b$ khác nhau (theo thứ tự đó) là $7$, $6$.\\ 
				 Trường hợp này có $ 7 \cdot 6=42$ số.
				\item Nếu $d = 5$ ta xét số $\overline{ab5}$, với $a,b,c \in \{1;2;3;4;5;6;7\}$,  $a \neq 0$.\\
				Số cách chọn $a,b$ khác nhau (theo  thứ tự đó) là $6$, $6$.\\
				Trường hợp này có $6 \cdot 6=36$  số.
			\end{itemize}
			Vậy có $42+36=78$ số thỏa mãn  yêu cầu của bài toán.
			\itemch Gọi số cần  tìm có dạng $\overline{abc}$, với $a$, $b$, $c$ khác nhau lấy từ $A=\{0;1;2;3;4;5;6;7\}$ và $a \neq 0$.
			\begin{itemize}
				\item Chọn $a<4$: có $3$ cách; chọn $b$: có $7$ cách; chọn $c$: có $6$ cách.\\ 
				Vậy theo quy tắc nhân có $  7 \cdot 6=42$ số.
				 \item Chọn $a=4$; chọn $b<5 \Rightarrow b \in \{0;1;2;3\}$: có $4$ cách; chọn $c$: có $6$ cách.\\
				Vậy theo quy tắc nhân có  $ 4 \cdot 6=24$ số.
				\item Chọn $a=4$; chọn $b=5$; chọn $c<6 \Rightarrow c  \in \{0;1;2;3\}$: có $4$ cách.\\
				Vậy  có $4$ số.
			\end{itemize}
			Vậy có $126+24+4=154$ số thỏa mãn yêu cầu của bài toán. 
		\end{itemchoice}
	 }
\end{ex}

% nguồn: MaHoaNH_DungSai_Lop10
\begin{ex}%[0D8C1-5]%[ma:caf07c03]
	Cho tập hợp $X=\{1;2;3;4;5;6;7;8;9\}$. 
	\choiceTF 
		{Có $3024$  số tự nhiên có $4$ chữ số lấy từ tập $X$}
		{\True Có $15\,120$ số tự nhiên có $5$ chữ số đôi một khác  nhau lấy từ tập $X$}
		{\True Có $1\,680$ số tự nhiên lẻ có $4$ chữ số đôi một khác nhau lấy từ tập  $X$}
		{Có $2\,520$ số tự nhiên có $4$ chữ số đôi một khác nhau lấy từ tập $X$ sao cho không có hai chữ số  liên tiếp nào cùng chẵn}
	 \loigiai{
		\begin{itemchoice}
			\itemch Gọi số tự nhiên có $4$ chữ số cần tìm là $\overline{abcd}$ ($a\neq0$; $0\leq a,b,c,d\leq  9$; $a,b,c,d\in\mathbb{N}$).\\
			\textit{Chọn số  tự nhiên có $4$ chữ số lấy từ tập $X$:}\\
			Chọn $a$ có $9$ cách, chọn $b$ có  $9$ cách, chọn $c$ có $9$ cách, chọn $d$ có $9$ cách.\\
			Theo quy  tắc nhân có $9\cdot9\cdot9\cdot9=6\,561$ số thỏa mãn.
			\itemch Gọi số tự nhiên có $5$ chữ số khác nhau cần tìm là $\overline{abcde}$ ($a\neq0$; $0\leq a,b,c,d,e\leq 9$;  $a,b,c,d,e\in\mathbb{N}$).\\
			Chọn $a$, $b$, $c$, $d$, $e$ có  $\mathrm{A}_9^5$ cách.\\
			Vậy  có $\mathrm{A}_{9}^{5}=15\,120$ số thỏa mãn.
			\itemch \textit{Chọn  số tự nhiên lẻ có $4$ chữ số lấy từ tập $X$:}\\
			Gọi số tự nhiên lẻ có $4$ chữ số cần tìm  là $\overline{abcd}$ ($a\neq0$; $0\leq a,b,c,d\leq 9$; $a,b,c,d\in\mathbb{N}$).\\
			Chọn $d$ có $5$ cách, chọn $a$, $b$,  $c$ có $\mathrm{A}_{8}^{3}$ cách.\\
			Theo quy tắc nhân  có $5\cdot\mathrm{A}_{8}^{3}=1\,680$ số thỏa mãn.
			\itemch \textit{Chọn số tự nhiên có $4$ chữ  số đôi một khác nhau lấy từ tập $X$ sao cho không có hai chữ số liên tiếp nào cùng chẵn:}\\
			Gọi số tự nhiên  có $4$ chữ số cần tìm là $\overline{abcd}$ ($a\neq0$; $0\leq a,b,c,d\leq 9$; $a,b,c,d\in\mathbb{N}$).
			\begin{itemize}
				\item \textbf{Trường hợp 1:} Bốn chữ số đều là số lẻ.\\ 
				Chọn $4$ số lẻ từ $5$ số lẻ và sắp xếp có  $\mathrm{A}_{5}^{4}=120$ cách.\\
				Do đó trường hợp này  có $120$ số thỏa mãn.
				 \item \textbf{Trường hợp 2:} Ba chữ số lẻ và một chữ số chẵn.\\
				Chọn $3$ chữ số lẻ từ $5$ chữ số lẻ có  $\mathrm{C}_{5}^{3}$ cách.\\
				Chọn $1$ chữ số chẵn từ $4$ chữ số chẵn có  $\mathrm{C}_{4}^{1}$ cách.\\
				Sắp xếp $4$ chữ số đã chọn  có $4!$ cách.\\
				 Theo quy tắc nhân có $\mathrm{C}_{5}^{3}\cdot\mathrm{C}_{4}^{1}\cdot4!=960$ số thỏa mãn.
				\item \textbf{Trường hợp 3:} Hai  chữ số lẻ và hai chữ số chẵn.\\
				Chọn $2$ chữ số lẻ từ  $5$ chữ số lẻ có $\mathrm{C}_{5}^{2}$ cách.\\
				 Chọn $2$ chữ số chẵn từ $4$ chữ số chẵn có $\mathrm{C}_{4}^{2}$ cách.\\
				Xếp $2$ chữ số lẻ có $2!$ cách.\\ 
				Xếp $2$ chữ số chẵn vào  $2$ trong $3$ vị trí (ở $2$ đầu hoặc ở giữa hai số lẻ) có $\mathrm{A}_{3}^{2}$ cách.\\
				Theo quy  tắc nhân có $\mathrm{C}_{5}^{2}\cdot\mathrm{C}_{4}^{2}\cdot2!\cdot\mathrm{A}_{3}^{2}=720$ số thỏa mãn.
			\end{itemize}
			Kết hợp ba trường hợp, theo quy tắc cộng  có $120+960+720=1\,800$ số thỏa mãn.
		\end{itemchoice}
	 }
\end{ex}


%%%%%%%%%%%%%%%% TRẢ LỜI NGẮN %%%%%%%%%%%%%%%%

% nguồn: MaHoaNH_TraLoiNgan_Lop10
\begin{bt}%[0D8N1-5]%[ma:50ef62c2]
	Từ 6 chữ số $1$, $2$, $3$, $4$, $5$, $6$ có thể  lập được bao nhiêu số có $2$ chữ số khác nhau và chia hết cho $9$?
	\shortans{4} 
	\loigiai{ 
		Số có hai chữ số chia  hết cho $9$ nếu tổng của chúng chia hết cho $9$, trong $6$ chữ số chỉ có $2$ cặp số $\{4,5\}$ và $\{3,6\}$. Do đó ta có $4$ số thỏa mãn yêu cầu là $45$, $54$, $36$, $63$.
	 }
\end{bt}

% nguồn: MaHoaNH_TraLoiNgan_Lop10
\begin{bt}%[0D8H1-5]%[ma:200d231c]
	Cho tập hợp $A=\{1;2;3;4;5;6;7;8\}$. Từ tập hợp $A$ có thể lập được bao nhiêu số gồm $8$ chữ số đôi một khác nhau sao cho các số này lẻ và không chia hết cho $5$. Khi đó, tổng các chữ số của kết quả  bằng bao nhiêu?
	 \shortans{9}
	 \loigiai{
		\begin{itemize}
			\item Chữ số cuối có $3$ cách chọn là $\{1;3;7\}$. 
			\item Số cách chọn các chữ số còn lại là $7 \cdot 6 \cdot 5  \cdot 4 \cdot 3 \cdot 2 \cdot 1 \Rightarrow 15120$ số cần tìm.
		\end{itemize}
	} 
\end{bt}

% nguồn: MaHoaNH_TraLoiNgan_Lop10
\begin{bt}%[0D8H1-5]%[ma:a25576fd]
	Từ các chữ số $0$, $1$, $2$, $3$, $4$, $5$ có thể lập được bao nhiêu số chẵn gồm $4$  chữ số khác nhau?
	\shortans{196} 
	 \loigiai{
		Gọi số chẵn có  bốn chữ số khác nhau có dạng là $\overline{a b c d}$ với $d \in\{0,2,4\}$.
		\begin{enumerate}[1.]
			\item  Trường hợp 1: Chọn $d=0$ có $1$ cách. Khi đó
			\begin{itemize}
				\item Chọn $a$ có $5$ cách vì $a  \neq 0$.
				\item Chọn $b$ có  $5$ cách vì $b \neq a$.
				 \item Chọn $c$ có $4$ cách vì $c \neq a$, $c \neq b$.
			\end{itemize}
			Theo quy tắc nhân số các  số là $1 \cdot 5 \cdot 5 \cdot 4=100$ số.
			\item Trường hợp 2: Chọn  $d \in\{2,4\}$ có $2$ cách. Khi đó
			\begin{itemize}
				\item Chọn $a$ có $4$ cách vì $a \neq 0$,  $a \neq d$.
				\item Chọn $b$ có $4$ cách vì  $b \neq a$, $b \neq d$.
				\item Chọn $c$ có $3$ cách vì $c \neq a$, $c \neq d$, $c \neq b$. 
			\end{itemize}
			Theo quy  tắc nhân số các số là $2 \cdot 4 \cdot 4 \cdot 3=96$ số.
		\end{enumerate}
		Vậy theo quy tắc cộng số các số cần tìm là $100+96=196$ số. 
	} 
\end{bt}
